% Options for packages loaded elsewhere
% Options for packages loaded elsewhere
\PassOptionsToPackage{unicode}{hyperref}
\PassOptionsToPackage{hyphens}{url}
\PassOptionsToPackage{dvipsnames,svgnames,x11names}{xcolor}
%
\documentclass[
]{agujournal2019}
\usepackage{xcolor}
\usepackage{amsmath,amssymb}
\setcounter{secnumdepth}{5}
\usepackage{iftex}
\ifPDFTeX
  \usepackage[T1]{fontenc}
  \usepackage[utf8]{inputenc}
  \usepackage{textcomp} % provide euro and other symbols
\else % if luatex or xetex
  \usepackage{unicode-math} % this also loads fontspec
  \defaultfontfeatures{Scale=MatchLowercase}
  \defaultfontfeatures[\rmfamily]{Ligatures=TeX,Scale=1}
\fi
\usepackage{lmodern}
\ifPDFTeX\else
  % xetex/luatex font selection
\fi
% Use upquote if available, for straight quotes in verbatim environments
\IfFileExists{upquote.sty}{\usepackage{upquote}}{}
\IfFileExists{microtype.sty}{% use microtype if available
  \usepackage[]{microtype}
  \UseMicrotypeSet[protrusion]{basicmath} % disable protrusion for tt fonts
}{}
\makeatletter
\@ifundefined{KOMAClassName}{% if non-KOMA class
  \IfFileExists{parskip.sty}{%
    \usepackage{parskip}
  }{% else
    \setlength{\parindent}{0pt}
    \setlength{\parskip}{6pt plus 2pt minus 1pt}}
}{% if KOMA class
  \KOMAoptions{parskip=half}}
\makeatother
% Make \paragraph and \subparagraph free-standing
\makeatletter
\ifx\paragraph\undefined\else
  \let\oldparagraph\paragraph
  \renewcommand{\paragraph}{
    \@ifstar
      \xxxParagraphStar
      \xxxParagraphNoStar
  }
  \newcommand{\xxxParagraphStar}[1]{\oldparagraph*{#1}\mbox{}}
  \newcommand{\xxxParagraphNoStar}[1]{\oldparagraph{#1}\mbox{}}
\fi
\ifx\subparagraph\undefined\else
  \let\oldsubparagraph\subparagraph
  \renewcommand{\subparagraph}{
    \@ifstar
      \xxxSubParagraphStar
      \xxxSubParagraphNoStar
  }
  \newcommand{\xxxSubParagraphStar}[1]{\oldsubparagraph*{#1}\mbox{}}
  \newcommand{\xxxSubParagraphNoStar}[1]{\oldsubparagraph{#1}\mbox{}}
\fi
\makeatother


\usepackage{longtable,booktabs,array}
\usepackage{calc} % for calculating minipage widths
% Correct order of tables after \paragraph or \subparagraph
\usepackage{etoolbox}
\makeatletter
\patchcmd\longtable{\par}{\if@noskipsec\mbox{}\fi\par}{}{}
\makeatother
% Allow footnotes in longtable head/foot
\IfFileExists{footnotehyper.sty}{\usepackage{footnotehyper}}{\usepackage{footnote}}
\makesavenoteenv{longtable}
\usepackage{graphicx}
\makeatletter
\newsavebox\pandoc@box
\newcommand*\pandocbounded[1]{% scales image to fit in text height/width
  \sbox\pandoc@box{#1}%
  \Gscale@div\@tempa{\textheight}{\dimexpr\ht\pandoc@box+\dp\pandoc@box\relax}%
  \Gscale@div\@tempb{\linewidth}{\wd\pandoc@box}%
  \ifdim\@tempb\p@<\@tempa\p@\let\@tempa\@tempb\fi% select the smaller of both
  \ifdim\@tempa\p@<\p@\scalebox{\@tempa}{\usebox\pandoc@box}%
  \else\usebox{\pandoc@box}%
  \fi%
}
% Set default figure placement to htbp
\def\fps@figure{htbp}
\makeatother


% definitions for citeproc citations
\NewDocumentCommand\citeproctext{}{}
\NewDocumentCommand\citeproc{mm}{%
  \begingroup\def\citeproctext{#2}\cite{#1}\endgroup}
\makeatletter
 % allow citations to break across lines
 \let\@cite@ofmt\@firstofone
 % avoid brackets around text for \cite:
 \def\@biblabel#1{}
 \def\@cite#1#2{{#1\if@tempswa , #2\fi}}
\makeatother
\newlength{\cslhangindent}
\setlength{\cslhangindent}{1.5em}
\newlength{\csllabelwidth}
\setlength{\csllabelwidth}{3em}
\newenvironment{CSLReferences}[2] % #1 hanging-indent, #2 entry-spacing
 {\begin{list}{}{%
  \setlength{\itemindent}{0pt}
  \setlength{\leftmargin}{0pt}
  \setlength{\parsep}{0pt}
  % turn on hanging indent if param 1 is 1
  \ifodd #1
   \setlength{\leftmargin}{\cslhangindent}
   \setlength{\itemindent}{-1\cslhangindent}
  \fi
  % set entry spacing
  \setlength{\itemsep}{#2\baselineskip}}}
 {\end{list}}
\usepackage{calc}
\newcommand{\CSLBlock}[1]{\hfill\break\parbox[t]{\linewidth}{\strut\ignorespaces#1\strut}}
\newcommand{\CSLLeftMargin}[1]{\parbox[t]{\csllabelwidth}{\strut#1\strut}}
\newcommand{\CSLRightInline}[1]{\parbox[t]{\linewidth - \csllabelwidth}{\strut#1\strut}}
\newcommand{\CSLIndent}[1]{\hspace{\cslhangindent}#1}



\setlength{\emergencystretch}{3em} % prevent overfull lines

\providecommand{\tightlist}{%
  \setlength{\itemsep}{0pt}\setlength{\parskip}{0pt}}



 


\usepackage{url} %this package should fix any errors with URLs in refs.
\usepackage{lineno}
\usepackage[inline]{trackchanges} %for better track changes. finalnew option will compile document with changes incorporated.
\usepackage{soul}
\linenumbers
\makeatletter
\@ifpackageloaded{caption}{}{\usepackage{caption}}
\AtBeginDocument{%
\ifdefined\contentsname
  \renewcommand*\contentsname{Table of contents}
\else
  \newcommand\contentsname{Table of contents}
\fi
\ifdefined\listfigurename
  \renewcommand*\listfigurename{List of Figures}
\else
  \newcommand\listfigurename{List of Figures}
\fi
\ifdefined\listtablename
  \renewcommand*\listtablename{List of Tables}
\else
  \newcommand\listtablename{List of Tables}
\fi
\ifdefined\figurename
  \renewcommand*\figurename{Figure}
\else
  \newcommand\figurename{Figure}
\fi
\ifdefined\tablename
  \renewcommand*\tablename{Table}
\else
  \newcommand\tablename{Table}
\fi
}
\@ifpackageloaded{float}{}{\usepackage{float}}
\floatstyle{ruled}
\@ifundefined{c@chapter}{\newfloat{codelisting}{h}{lop}}{\newfloat{codelisting}{h}{lop}[chapter]}
\floatname{codelisting}{Listing}
\newcommand*\listoflistings{\listof{codelisting}{List of Listings}}
\makeatother
\makeatletter
\makeatother
\makeatletter
\@ifpackageloaded{caption}{}{\usepackage{caption}}
\@ifpackageloaded{subcaption}{}{\usepackage{subcaption}}
\makeatother
\usepackage{bookmark}
\IfFileExists{xurl.sty}{\usepackage{xurl}}{} % add URL line breaks if available
\urlstyle{same}
\hypersetup{
  pdftitle={The consequences of climate change in Atlantic Forest birds},
  pdfauthor={Eduardo S.A. Santos; João C.T. Menezes; Gustavo Burin},
  pdfkeywords={body size, climate change, morphological
adaptation, biotic responses, trophic-interaction},
  colorlinks=true,
  linkcolor={blue},
  filecolor={Maroon},
  citecolor={Blue},
  urlcolor={Blue},
  pdfcreator={LaTeX via pandoc}}


\journalname{OSF Registries \textbar{} Open-Ended Registration}

\draftfalse

\begin{document}
\title{The consequences of climate change in Atlantic Forest birds}

\authors{Eduardo S.A. Santos\affil{1}, João C.T.
Menezes\affil{2}, Gustavo Burin\affil{3}}
\affiliation{1}{Collaboration for Open Science and Synthesis in Ecology
and Evolution, Department of Biological Sciences, University of Alberta,
Edmonton, AB, T6G 2E9, Canada, }\affiliation{2}{Organismic and
Evolutionary Biology Graduate Program, University of Massachusetts
Amherst, Amherst, MA, 01003, USA, }\affiliation{3}{Department of
Biological and Environmental Sciences, University of Gothenburg,
Sweden, }
\correspondingauthor{Eduardo S.A. Santos}{esantos2@ualberta.ca}


\begin{abstract}
Here we present a preregistration for a phylogenetic comparative study
of the effects of climate change on morphology of Atlantic Forest birds.
Below we will provide detailed information on the planned research
project.
\end{abstract}

\section*{Plain Language Summary}
Earthquake data for the island of La Palma from the September 2021
eruption is found \ldots{}




\section{Introduction}\label{introduction}

A well-known ecogeographical pattern describes the relationship between
environmental temperature variation and body size. This association,
known as Bergmann's rule (Bergmann, 1847), states that animals in warmer
climates tend to be smaller in size. The logic underlying this rule is
that the increase in body surface area to volume ratio should facilitate
more efficient thermoregulation (Mayr, 1956). Researchers have found
empirical evidence supporting Bergmann's rule both within and across
species, specifically for birds (Ashton, 2002; Frӧhlich et al., 2023;
McQueen et al., 2022). The relationship between warmer temperatures and
smaller body sizes has led researchers for decades to predict that
anthropogenic climate change may cause changes in body size in a
analogous way as the ecogeographical pattern. That is, as the Earth's
climate warms, body size in endothermic animals, such as birds, should
decrease (Gardner et al., 2011; Sheridan \& Bickford, 2011; Celine
Teplitsky \& Millien, 2014).

Two of the main responses of organisms to climate change are
phenological changes and geographical range shifts (Parmesan \& Yohe,
2003). Yet, the shrinkage of bodies, a third response of organisms to
global warming, is increasingly being observed (Gardner et al., 2011;
Sheridan \& Bickford, 2011; Celine Teplitsky \& Millien, 2014). Body
size shrinkage due to climate change has been hypothesized to be a
consequence of ecological or evolutionary processes (Ozgul et al., 2009;
Ryding et al., 2024; Weeks et al., 2020).

The ecological hypothesis posits that climate change influences trophic
interactions, thus leading to a reduction in body size.

Changes in temperature can result in temperature-dependent or
-independent mechanisms that influence body size (reviewed in Gardner et
al., 2011). Among temperature-independent mechanisms, decreases in food
availability caused by habitat loss or fragmentation may lead to the
reduction of body size (Schmidt \& Jensen, 2005). Thus, limiting
juveniles access to resources could lead to smaller adult body size
(Rode et al., 2010; Céline Teplitsky et al., 2008). Reductions in body
size have been attributed to genetic microevolutionary responses to
warming; smaller individuals dissipate heat better than larger
individuals due to the larger ratio between surface:volume of their
bodies (\emph{e.g.}, Bergmann's rule, Celine Teplitsky \& Millien,
2014).\footnote{I think that this paragraph is not necessary. Talk only
  about the ecological consequences.}

\subsection{Research Question}\label{research-question}

Here, we will conduct a large scale comparative phylogenetic analysis of
Neotropical passerines that inhabit the Brazilian Atlantic Forest to
investigate whether temporal changes in climate variables are associated
with body size changes.

\subsection{Hypothesis \& Predictions}\label{hypothesis-predictions}

\section{Data \& Methods}\label{sec-data-methods}

\subsection{Data Collection}\label{data-collection}

\subsection{Target Sample Sizes}\label{target-sample-sizes}

\subsection{Data Analysis Plan}\label{data-analysis-plan}

\section{Conclusion}\label{conclusion}

\section*{References}\label{references}
\addcontentsline{toc}{section}{References}

\phantomsection\label{refs}
\begin{CSLReferences}{1}{0}
\vspace{1em}

\bibitem[\citeproctext]{ref-ashton2002a}
Ashton, K. G. (2002). Patterns of within-species body size variation of
birds: strong evidence for bergmann's rule. \emph{Global Ecology and
Biogeography}, \emph{11}(6), 505--523.
\url{https://doi.org/10.1046/j.1466-822X.2002.00313.x}

\bibitem[\citeproctext]{ref-bergmann1847}
Bergmann, C. (1847). Über die verhältnisse der warmeokonomiedie thiere
zu ihrer grösse. \emph{Göttinger Studien}, \emph{3}, 359--708. Retrieved
from \url{https://www.digitale-sammlungen.de/en/view/bsb10306637?page=1}

\bibitem[\citeproctext]{ref-frohlich2023a}
Frӧhlich, A., Kotowska, D., Martyka, R., \& Symonds, M. R. E. (2023).
Allometry reveals trade-offs between bergmann{'}s and allen{'}s rules,
and different avian adaptive strategies for thermoregulation.
\emph{Nature Communications}, \emph{14}(1), 1101.
\url{https://doi.org/10.1038/s41467-023-36676-w}

\bibitem[\citeproctext]{ref-gardner_declining_2011}
Gardner, J. L., Peters, A., Kearney, M. R., Joseph, L., \& Heinsohn, R.
(2011). Declining body size: A third universal response to warming?
\emph{Trends in Ecology \& Evolution}, \emph{26}(6), 285--291.
\url{https://doi.org/10.1016/j.tree.2011.03.005}

\bibitem[\citeproctext]{ref-mayr1956}
Mayr, E. (1956). Geographical character gradients and climatic
adaptation. \emph{Evolution}, \emph{10}(1), 105--108.
\url{https://doi.org/10.2307/2406103}

\bibitem[\citeproctext]{ref-mcqueen2022a}
McQueen, A., Klaassen, M., Tattersall, G. J., Atkinson, R., Jessop, R.,
Hassell, C. J., et al. (2022). Thermal adaptation best explains
bergmann{'}s and allen{'}s rules across ecologically diverse shorebirds.
\emph{Nature Communications}, \emph{13}(1), 4727.
\url{https://doi.org/10.1038/s41467-022-32108-3}

\bibitem[\citeproctext]{ref-ozgul2009}
Ozgul, A., Tuljapurkar, S., Benton, T. G., Pemberton, J. M.,
Clutton-Brock, T. H., \& Coulson, T. (2009). The Dynamics of Phenotypic
Change and the Shrinking Sheep of St. Kilda. \emph{Science},
\emph{325}(5939), 464--467. \url{https://doi.org/10/db5kwc}

\bibitem[\citeproctext]{ref-parmesan_globally_2003}
Parmesan, C., \& Yohe, G. (2003). A globally coherent fingerprint of
climate change impacts across natural systems. \emph{Nature},
\emph{421}(6918), 37. \url{https://doi.org/10.1038/nature01286}

\bibitem[\citeproctext]{ref-rode_reduced_2010}
Rode, K. D., Amstrup, S. C., \& Regehr, E. V. (2010). Reduced body size
and cub recruitment in polar bears associated with sea ice decline.
\emph{Ecological Applications}, \emph{20}(3), 768--782.
\url{https://doi.org/10.1890/08-1036.1}

\bibitem[\citeproctext]{ref-ryding2024}
Ryding, S., McQueen, A., Klaassen, M., Tattersall, G. J., \& Symonds, M.
R. E. (2024). Long- and short-term responses to climate change in body
and appendage size of diverse Australian birds. \emph{Global Change
Biology}, \emph{30}(10), e17517. \url{https://doi.org/10.1111/gcb.17517}

\bibitem[\citeproctext]{ref-schmidt_concomitant_2005}
Schmidt, N., \& Jensen, P. (2005). Concomitant {Patterns} in {Avian} and
{Mammalian} {Body} {Length} {Changes} in {Denmark}. \emph{Ecology and
Society}, \emph{10}(2). \url{https://doi.org/10/gf389w}

\bibitem[\citeproctext]{ref-sheridan_shrinking_2011}
Sheridan, J. A., \& Bickford, D. (2011). Shrinking body size as an
ecological response to climate change. \emph{Nature Climate Change},
\emph{1}(8), 401--406. \url{https://doi.org/10.1038/nclimate1259}

\bibitem[\citeproctext]{ref-teplitsky_climate_2014}
Teplitsky, Celine, \& Millien, V. (2014). Climate warming and
{Bergmann}'s rule through time: Is there any evidence?
\emph{Evolutionary Applications}, \emph{7}(1), 156--168.
\url{https://doi.org/10.1111/eva.12129}

\bibitem[\citeproctext]{ref-teplitsky_bergmanns_2008}
Teplitsky, Céline, Mills, J. A., Alho, J. S., Yarrall, J. W., \& Merilä,
J. (2008). Bergmann's rule and climate change revisited: {Disentangling}
environmental and genetic responses in a wild bird population.
\emph{Proceedings of the National Academy of Sciences}, \emph{105}(36),
13492--13496. \url{https://doi.org/10.1073/pnas.0800999105}

\bibitem[\citeproctext]{ref-weeks2020}
Weeks, B. C., Willard, D. E., Zimova, M., Ellis, A. A., Witynski, M. L.,
Hennen, M., \& Winger, B. M. (2020). Shared morphological consequences
of global warming in North American migratory birds. \emph{Ecology
Letters}, \emph{23}(2), 316--325.
\url{https://doi.org/10.1111/ele.13434}

\end{CSLReferences}




\end{document}
